\clearpage
\VerdierPhysicalPageAnchor{0125}%
\begingroup
\small
\setlength{\emergencystretch}{1.5em}%
\setlength{\abovedisplayskip}{4pt plus 2pt minus 2pt}%
\setlength{\belowdisplayskip}{4pt plus 2pt minus 2pt}%

\VerdierSemanticAnchor{ch2-1-3-5}
\noindent\textbf{Proposition 1.3.5.}\quad
\emph{Soit $C$ une catégorie abélienne. Le foncteur :}

\VerdierSemanticAnchor{ch2-1-3-5-h0-functor}
\[
  H^0\colon K(C)\longrightarrow C
\]
\noindent\emph{est un foncteur cohomologique.}

\VerdierSemanticAnchor{ch2-1-3-5-first-exact-sequence}
En effet, d'après (chap. I, 3.3.1), il suffit de montrer que, pour tout objet
$X\xrightarrow{f}Y$ de $\operatorname{tr}_1(C)$ (chap. I, 3.1.1), la suite :
\[
  H^0(X)\xrightarrow{H^0(f)}H^0(Y)
  \xrightarrow{H^0(p(f))}H^0(c(f))
  \xrightarrow{H^0(q(f))}H^0(X[1])
\]
\noindent est exacte.

\VerdierSemanticAnchor{ch2-1-3-5-trii-reduction}
En utilisant la propriété TRII (1.1.1), on voit qu'il suffit de montrer que
la suite :

\VerdierSemanticAnchor{ch2-1-3-5-reduced-exact-sequence}
\[
  H^0(Y)\xrightarrow{H^0(p(f))}H^0(c(f))
  \xrightarrow{H^0(q(f))}H^0(X[1])
\]
\noindent est exacte, ce qui résulte immédiatement des définitions
(chap. I, 3.1.2 et 3.2.2).

\VerdierSemanticAnchor{ch2-1-3-5-proof-complete}
Le foncteur $H^0\colon K(C)\longrightarrow C$ est appelé le
\emph{foncteur cohomologique canonique}.

\VerdierSemanticAnchor{ch2-1-3-5-shifted-cohomology}
On pose $H^n(X)=H^0(X[n])$ conformément à l'usage.

\VerdierSemanticAnchor{ch2-1-3-6}
\medskip
\noindent\textbf{Proposition 1.3.6.}\quad
\emph{Soit $C$ une catégorie abélienne. Les conditions suivantes sont
équivalentes :}

\VerdierSemanticAnchor{ch2-1-3-6-equivalent-conditions}
\begin{enumerate}
\renewcommand{\labelenumi}{\roman{enumi})}
\item \emph{La catégorie $K(C)$ est abélienne.}
\item \emph{La catégorie $K(C)$ est décomposable (1.2.8).}
\item \emph{La catégorie $C$ est décomposable.}
\end{enumerate}

\VerdierSemanticAnchor{ch2-1-3-6-i-iff-ii}
L'équivalence
$\mathrm{i)}\Longleftrightarrow\mathrm{ii)}$ résulte de (1.2.9).

\VerdierSemanticAnchor{ch2-1-3-6-ii-implies-iii}
L'implication $\mathrm{ii)}\Longrightarrow\mathrm{iii)}$ est évidente car
$C$ est une sous-catégorie pleine de $K(C)$.

\VerdierSemanticAnchor{ch2-1-3-6-hstar-complex}
Pour tout complexe $X$ de $C$, désignons par $H^*(X)$ le complexe :
\[
  H^*(X)^l=H^l(X),\qquad l\in\mathbb Z,
\]
\[
  d_{H^*(X)}^l=0,\qquad l\in\mathbb Z.
\]

\VerdierSemanticAnchor{ch2-1-3-6-hstar-equivalence}
On vérifie que le foncteur
$H^*\colon K(C)\longrightarrow K(C)$ est une équivalence de catégories
lorsque $C$ est décomposable.

\VerdierSemanticAnchor{ch2-1-3-6-proof-complete}
On en déduit l'implication
$\mathrm{iii)}\Longrightarrow\mathrm{ii)}$.
\endgroup
